% Chapter 2, section 7: Advanced transactional usage
% Write in plain LaTeX. Code of a guide step: \guide{<step>}; commands: \begin{shell}...\end{shell};
% a screenshot: \screenshot{file}{caption}; a gap in the draft: \placeholder{...}.
% See book/latex/queen.sty for the full list.

\section{Advanced transactional usage}

The ack, the order's new state, the next event and tomorrow's reminder commit as one log entry, or
none of them does. That one call replaces the outbox table, the idempotency table and the cron job.

\guide{transaction}

Then open the dashboard: its KV page holds the order's new state, and its Timers page holds
tomorrow's reminder.

\why{a worker's step is several writes: it marks its input done, saves what changed, tells the
next service and plans what comes later. When those writes go to different systems, a crash between
two of them leaves the step half done, and someone writes code to find and repair the halves. One
commit has no halves.}

If the worker took too long and its lease ran out, the whole transaction is refused, so a slow
worker can never write over the one that has the message now. If the step committed and only the
answer got lost, a retry with the same \texttt{transactionId}s is refused whole, and nothing is
written twice.

\queenfig{lease-fence}{A worker whose lease ran out: its whole transaction rolls back.}

On the site, with the code in all six SDK languages: \href{https://queenmq.com/writing-applications/transactions/}{queenmq.com/writing-applications/transactions}.
